% questions.tex —— 1990 年·广东卷·理科 · 题目内容片段（无答案）
% 由 tools/generate_projects.py 依来源总账自 1990/1990_数学真题_广东理_试题.md 生成，勿手改（改源 MD 后重跑）


\gkpart{第一卷}

每题 2 分；不选、选错或者选出的代号超过一个，均得 0 分。

\begin{question}
\gkti{1}如果直线 $ax+2y+1=0$ 与直线 $x+y-2=0$ 互相垂直，那么 $a$ 的值等于（　）。

\begin{gkopts}(4)
\task A．$1$
\task B．$-\dfrac{1}{3}$
\task C．$-\dfrac{2}{3}$
\task D．$-2$
\end{gkopts}

\end{question}
\begin{question}
\gkti{2}如果直线 $l$ 是平面 $\alpha$ 的斜线，那么在平面 $\alpha$ 内（　）。

\begin{gkopts}(1)
\task A．不存在与 $l$ 平行的直线；B．不存在与 $l$ 垂直的直线；
\task C．与 $l$ 垂直的直线只有一条；D．与 $l$ 平行的直线有无穷多条。
\end{gkopts}

\end{question}
\begin{question}
\gkti{3}某乒乓球队有 $9$ 名队员，其中 $2$ 名是种子选手，现要挑选 $5$ 名队员参加比赛，种子选手都必须在内，那么不同的选法共有（　）。

\begin{gkopts}(4)
\task A．$126$ 种
\task B．$84$ 种
\task C．$35$ 种
\task D．$21$ 种
\end{gkopts}

\end{question}
\begin{question}
\gkti{4}把 $6$ 个不同的元素排成前后两排，每排 $3$ 个元素，那么不同的排法共有（　）。

\begin{gkopts}(4)
\task A．$36$ 种
\task B．$120$ 种
\task C．$720$ 种
\task D．$1440$ 种
\end{gkopts}

\end{question}
\begin{question}
\gkti{5}一个圆台的母线长是上、下底面半径的等差中项，且侧面积为 $8\pi\,\mathrm{cm}^2$，那么母线长为（　）。

\begin{gkopts}(2)
\task A．$4\,\mathrm{cm}$
\task B．$2\sqrt2\,\mathrm{cm}$
\task C．$2\,\mathrm{cm}$
\task D．$\sqrt2\,\mathrm{cm}$
\end{gkopts}

\end{question}
\begin{question}
\gkti{6}如果复数 $z=3+ai$ 满足条件 $|z-2|<2$，那么实数 $a$ 的取值范围是（　）。

\begin{gkopts}(2)
\task A．$(-2\sqrt2,2\sqrt2)$
\task B．$(-2,2)$
\task C．$(-1,1)$
\task D．$(-\sqrt3,\sqrt3)$
\end{gkopts}

\end{question}
\begin{question}
\gkti{7}一动点在圆 $x^2+y^2=1$ 上移动时，它与定点 $(3,0)$ 连线的中点的轨迹方程是（　）。

\begin{gkopts}(2)
\task A．$(x+3)^2+y^2=4$
\task B．$(x-3)^2+y^2=1$；
\task C．$(2x-3)^2+4y^2=1$
\task D．$\left(x+\dfrac32\right)^2+y^2=\dfrac12$。
\end{gkopts}

\end{question}
\begin{question}
\gkti{8}函数
\nopagebreak[4]
\[
y=\frac{\sin x}{|\sin x|}+\frac{|\cos x|}{\cos x}+\frac{\operatorname{tg}x}{|\operatorname{tg}x|}+\frac{|\operatorname{ctg}x|}{\operatorname{ctg}x}
\]
的值域是（　）。

\begin{gkopts}(4)
\task A．$\{-2,4\}$
\task B．$\{-2,0,4\}$
\task C．$\{-2,0,2,4\}$
\task D．$\{-4,-2,0,4\}$
\end{gkopts}

\end{question}
\begin{question}
\gkti{9}方程 $\sin 2x=\sin x$ 在区间 $(0,2\pi)$ 内的解的个数是（　）。

\begin{gkopts}(4)
\task A．$4$
\task B．$3$
\task C．$2$
\task D．$1$
\end{gkopts}

\end{question}
\begin{question}
\gkti{10}函数 $y=(x+4)^2$ 在某区间上是减函数，这区间可以是（　）。

\begin{gkopts}(4)
\task A．$(-\infty,-4]$
\task B．$[-4,+\infty)$
\task C．$[4,+\infty)$
\task D．$(-\infty,4]$
\end{gkopts}

\end{question}
\begin{question}
\gkti{11}已知函数：① $y=\operatorname{arctg}x$，② $y=\dfrac\pi2-\operatorname{arcctg}x$。那么（　）。

\begin{gkopts}(1)
\task A．①与②都是奇函数；B．①与②都是偶函数；
\task C．①是奇函数，②是偶函数；D．①与②都不是奇函数，也都不是偶函数。
\end{gkopts}

\end{question}
\begin{question}
\gkti{12}已知下图是函数 $y=2\sin(\omega x+\varphi)$（$|\varphi|<\dfrac\pi2$）的图象，那么（　）。

\begin{gktip}
\textbf{缺源提示：}固定来源 PDF 第 2 页该题图位为空白。以下只照录题干和选项；不依据答案表重建图形，题目保持待核。
\end{gktip}

\begin{gkopts}(1)
\task A．$\omega=2,\ \varphi=-\dfrac\pi6$
\task B．$\omega=2,\ \varphi=\dfrac\pi6$；
\task C．$\omega=\dfrac{10}{11},\ \varphi=\dfrac\pi6$
\task D．$\omega=\dfrac{10}{11},\ \varphi=-\dfrac\pi6$。
\end{gkopts}

\end{question}
\begin{question}
\gkti{13}如果直线 $y=ax+2$ 与直线 $y=3x-b$ 关于直线 $y=x$ 对称，那么（　）。

\begin{gkopts}(2)
\task A．$a=\dfrac13,\ b=-6$
\task B．$a=\dfrac13,\ b=6$；
\task C．$a=3,\ b=-2$
\task D．$a=3,\ b=6$。
\end{gkopts}

\end{question}
\begin{question}
\gkti{14}如果抛物线 $y^2=a(x+1)$ 的准线方程是 $x=-3$，那么它的焦点坐标是（　）。

\begin{gkopts}(4)
\task A．$(3,0)$
\task B．$(2,0)$
\task C．$(1,0)$
\task D．$(-1,0)$
\end{gkopts}

\end{question}
\begin{question}
\gkti{15}如果 $\log_a2>\log_b2>0$，那么（　）。

\begin{gkopts}(4)
\task A．$1<a<b$
\task B．$1<b<a$
\task C．$0<a<b<1$
\task D．$0<b<a<1$
\end{gkopts}

\end{question}
\begin{question}
\gkti{16}复数 $1+i\operatorname{ctg}\theta$（$\pi<\theta<2\pi$）的三角形式是（　）。

\begin{gkopts}(1)
\task A．$\dfrac1{\sin\theta}(\sin\theta+i\cos\theta)$；
\task B．$\dfrac1{\sin\theta}\left[\cos\left(\dfrac\pi2-\theta\right)+i\sin\left(\dfrac\pi2-\theta\right)\right]$；
\task C．$-\dfrac1{\sin\theta}\left[\cos\left(\dfrac\pi2+\theta\right)+i\sin\left(\dfrac\pi2+\theta\right)\right]$；
\task D．$-\dfrac1{\sin\theta}\left[\cos\left(\dfrac{3\pi}2-\theta\right)+i\sin\left(\dfrac{3\pi}2-\theta\right)\right]$。
\end{gkopts}

\end{question}
\begin{question}
\gkti{17}如果 $x,y$ 是实数，那么“$xy>0$”是“$|x+y|=|x|+|y|$”的（　）。

\begin{gkopts}(2)
\task A．充分但不必要条件；B．必要但不充分条件；
\task C．必要而且充分的条件；D．既不充分也不必要的条件。
\end{gkopts}

\end{question}
\begin{question}
\gkti{18}极坐标方程 $4\rho\sin^2\dfrac\theta2=5$ 表示的曲线是（　）。

\begin{gkopts}(4)
\task A．圆
\task B．椭圆
\task C．抛物线
\task D．双曲线的一支
\end{gkopts}

\end{question}
\begin{question}
\gkti{19}已知等比数列的公比为 $2$，且前 $4$ 项之和等于 $1$，那么前 $8$ 项之和等于（　）。

\begin{gkopts}(4)
\task A．$15$
\task B．$17$
\task C．$19$
\task D．$21$
\end{gkopts}

\end{question}
\begin{question}
\gkti{20}函数 $y=\operatorname{tg}\dfrac x2-\dfrac1{\sin x}$ 的最小正周期是（　）。

\begin{gkopts}(4)
\task A．$\dfrac\pi2$
\task B．$\pi$
\task C．$\dfrac{3\pi}2$
\task D．$2\pi$
\end{gkopts}

\end{question}
\begin{question}
\gkti{21}$\operatorname{tg}70^\circ+\operatorname{tg}50^\circ-\sqrt3\,\operatorname{tg}50^\circ\operatorname{tg}70^\circ$ 的值等于（　）。

\begin{gkopts}(2)
\task A．$\sqrt3$
\task B．$\dfrac{\sqrt3}3$
\task C．$-\dfrac{\sqrt3}3$
\task D．$-\sqrt3$
\end{gkopts}

\end{question}
\begin{question}
\gkti{22}如果实数 $x,y$ 满足等于 $(x-2)^2+y^2=3$，那么 $\dfrac yx$ 的最大值是（　）。

\begin{gkopts}(2)
\task A．$\dfrac12$
\task B．$\sqrt3$
\task C．$\dfrac{\sqrt3}2$
\task D．$\dfrac{\sqrt3}3$
\end{gkopts}

\end{question}
\begin{question}
\gkti{23}以一个正方体的顶点为顶点的四面体共有（　）。

\begin{gkopts}(4)
\task A．$70$ 个
\task B．$64$ 个
\task C．$58$ 个
\task D．$52$ 个
\end{gkopts}

\end{question}
\begin{question}
\gkti{24}设 $z$ 是复数，$z+2$ 的辐角为 $\dfrac\pi3$，$z-2$ 的辐角为 $\dfrac{5\pi}6$，那么 $z$ 等于（　）。

\begin{gkopts}(2)
\task A．$-\dfrac12+\dfrac{\sqrt3}2i$
\task B．$-\dfrac{\sqrt3}2+\dfrac12i$；
\task C．$-\sqrt3+i$
\task D．$-1+\sqrt3i$
\end{gkopts}

\end{question}
\begin{question}
\gkti{25}设函数 $y=\operatorname{arctg}x$ 的图象沿 $x$ 轴正方向平移 $2$ 个单位，所得图象为 $C$。又设图象 $C_1$ 与 $C$ 关于原点对称，那么 $C_1$ 所对应的函数是（　）。

\begin{gkopts}(2)
\task A．$y=-\operatorname{arctg}(x-2)$
\task B．$y=\operatorname{arctg}(x-2)$；
\task C．$y=-\operatorname{arctg}(x+2)$
\task D．$y=\operatorname{arctg}(x+2)$。
\end{gkopts}

\end{question}
\gkpart{第二卷}

\gkpart{一、填空题}

本题满分 15 分，共 5 小题，每小题 3 分。每题都有用横线表示的空位。

\Needspace{0.4\textheight}
\begin{question}
\gktiwithlabel{26}{一}

（1） 双曲线 $\dfrac{x^2}{16}-\dfrac{y^2}{9}=1$ 的准线方程是\gkblank{}。

（2） $(x-1)-(x-1)^2+(x-1)^3-(x-1)^4$ 的展开式中，$x^2$ 的系数等于\gkblank{}。

\Needspace{5\baselineskip}
（3） 已知数列 $\{a_n\}$ 是公差不为零的等差数列，如果 $S_n$ 是 $\{a_n\}$ 的前 $n$ 项的和，那么
\nopagebreak[4]
\[
\lim_{n\to\infty}\frac{na_n}{S_n}
\]
的值等于\gkblank{}。

（4） 函数 $y=\sin x+\cos x$ 的最大值是\gkblank{}。

（5） 如图，三棱柱 $ABC-A_1B_1C_1$ 中，若 $E,F$ 分别为 $AB,AC$ 的中点，平面 $EB_1C_1F$ 将三棱柱分成体积为 $V_1,V_2$ 的两部分，那么 $V_1:V_2$ 等于\gkblank{}。
\nopagebreak[4]
\begin{center}\begin{tikzpicture}[x=0.78cm,y=0.66cm,line cap=round,line join=round]
\coordinate (A) at (0,0);
\coordinate (B) at (4,0);
\coordinate (C) at (1.4,1.9);
\coordinate (Aone) at (-0.35,2.0);
\coordinate (Bone) at (3.65,2.0);
\coordinate (Cone) at (1.05,3.9);
\coordinate (E) at (2,0);
\coordinate (F) at (0.7,0.95);

% 原卷三棱柱及平面 EB_1C_1F 的拓扑关系；投影不表示长度比例。
\fill[gray!14] (E)--(Bone)--(Cone)--(F)--cycle;
\draw[thick] (A)--(B)--(C)--cycle;
\draw[thick] (Aone)--(Bone)--(Cone)--(Aone);
\draw[thick] (A)--(Aone) (B)--(Bone) (C)--(Cone);
\draw[dashed] (A)--(C);
\draw[very thick] (E)--(Bone) (Cone)--(F) (F)--(E) (Bone)--(Cone);

\node[below left] at (A) {$A$};
\node[below right] at (B) {$B$};
\node[above right] at (C) {$C$};
\node[above left] at (Aone) {$A_1$};
\node[above right] at (Bone) {$B_1$};
\node[above] at (Cone) {$C_1$};
\node[below] at (E) {$E$};
\node[left] at (F) {$F$};
\node at (0.3,1.55) {$V_1$};
\node at (3.4,1.0) {$V_2$};
\node[font=\small,align=center] at (1.7,-1.3)
  {依原卷关系重绘的示意图；非比例图。};
\end{tikzpicture}
\end{center}
\nopagebreak[4]
\end{question}
\gkpart{二、解答题}

\Needspace{0.4\textheight}
\begin{question}
\gktiwithlabel{27}{二}（本题满分 13 分）

如图，在三棱锥 $S-ABC$ 中，$SA\perp$ 底面 $ABC$，$AB\perp BC$，$DE$ 垂直平分 $SC$，且分别交 $AC,SC$ 于 $D,E$，又 $SA=AB=a$，$BC=\sqrt2a$。

（1） 求证：$SC\perp$ 面 $BDE$；

（2） 求面 $BDE$ 与面 $BDC$ 所成二面角的大小。
\nopagebreak[4]
\begin{center}\begin{tikzpicture}[x=0.95cm,y=0.78cm,line cap=round,line join=round]
\coordinate (A) at (-1.6,0);
\coordinate (C) at (2.4,0);
\coordinate (B) at (0.4,-1.65);
\coordinate (S) at (-1.6,3.1);
\coordinate (D) at (-0.8,0);
\coordinate (E) at (0.4,1.55);

% 依原卷中点、垂直平分线及棱锥的关系重绘；不表达额外长度比例。
\draw[thick] (A)--(S)--(C)--(B)--cycle;
\draw[thick] (S)--(B);
\draw[dashed] (A)--(C) (D)--(B);
\draw[dashed,thick] (D)--(E);
\draw[thick] (S)--(E);
\draw[thick] (E)--(C);

% SC 的等长刻痕。D、E 的投影位置保留原卷的垂直平分线关系。
\draw (-0.71,2.185) -- (-0.49,2.465);
\draw (1.29,0.635) -- (1.51,0.915);
% DE 与 SC 的直角标记，按两边方向构造小方角。
\draw (0.266,1.376) -- (0.44,1.242) -- (0.574,1.416);

\node[above left] at (S) {$S$};
\node[left] at (A) {$A$};
\node[below] at (B) {$B$};
\node[right] at (C) {$C$};
\node[below] at (D) {$D$};
\node[above right] at (E) {$E$};
\node[font=\small,align=center] at (0.4,-2.35)
  {依原卷关系重绘的示意图；非比例图。};
\end{tikzpicture}
\end{center}
\nopagebreak[4]
\end{question}
\begin{question}
\gktiwithlabel{28}{三}（本题满分 14 分）

已知椭圆的焦点为 $F_1(0,-1)$ 和 $F_2(0,1)$，直线 $y=4$ 是椭圆的一条准线。

（1） 求椭圆的方程；

（2） 又设点 $P$ 在这个椭圆上，且 $|PF_1|-|PF_2|=1$，求 $\operatorname{tg}\angle F_1PF_2$ 的值。

\end{question}
\begin{question}
\gktiwithlabel{29}{四}（本题满分 14 分）

\Needspace{5\baselineskip}
已知数列 $\{a_n\}$ 的前 $n$ 项和的公式是
\nopagebreak[4]
\[
S_n=\frac\pi{12}(2n^2+n)
\]
。

（1） 求证 $\{a_n\}$ 是等差数列，并求出它的首项和公差；

\Needspace{5\baselineskip}
（2） 记 $b_n=\sin a_n\cdot\sin a_{n+1}\cdot\sin a_{n+2}$，求证：对任何自然数 $n$，都有
\nopagebreak[4]
\[
b_n=\frac{\sqrt2}{8}(-1)^{n-1}
\]
。

\end{question}
\begin{question}
\gktiwithlabel{30}{五}（本题满分 14 分）

\Needspace{5\baselineskip}
设
\nopagebreak[4]
\[
f(x)=\lg\frac{1+2^x+4^xa}{3},\qquad a\in\mathbb R
\]
。

（1） 如果 $0<a\le1$，求证：当 $x\ne0$ 时，有 $2f(x)<f(2x)$；

（2） 如果当 $x\in(-\infty,1]$ 时 $f(x)$ 有意义，求 $a$ 的取值范围。

\begin{gktip}
完整试题共 30 个内部题目块。第 12 题图位缺失状态保持待核；本文件仅为来源转录，来源答案与独立解析另行登记。
\end{gktip}

\end{question}